% ID: SEL-D0067-OB-E03
% ESTADO: cotejada. Segunda pasada visual y borrador compilado revisados.
% Fuente: Selectividad/sel_2009_matematicas_aplicadas/mat apli CC SS01.pdf
% SHA-256: f5e34036805266809f1b5a2c6087d7ae598afd929b5179ec5ea3888db6cbccf4
% Página: 2; opción B; ejercicio original 3.

\noindent\underline{Parte I}

Una encuesta realizada por un banco muestra que el $60\%$ de sus clientes tiene un préstamo hipotecario, el $50\%$ tiene un préstamo personal y el $20\%$ tiene un préstamo de cada tipo. Se elige, al azar, un cliente de ese banco.
\begin{enumerate}[label=\alph*)]
\item \textbf{(1 punto)} Calcule la probabilidad de que no tenga ninguno de los dos préstamos.
\item \textbf{(1 punto)} Calcule la probabilidad de que tenga un préstamo hipotecario, sabiendo que no tiene un préstamo personal.
\end{enumerate}
\begin{samepage}
\noindent\underline{Parte II}

El cociente intelectual de los alumnos de un centro educativo se distribuye según una ley Normal de media 110 y desviación típica 15. Se extrae una muestra aleatoria simple de 25 alumnos.
\begin{enumerate}[label=\alph*)]
\item \textbf{(1.5 puntos)} \textquestiondown{}Cuál es la probabilidad de que la media del cociente intelectual de los alumnos de esa muestra sea superior a 113?
\item \textbf{(0.5 puntos)} Razone cómo se vería afectada la respuesta a la pregunta anterior si el tamaño de la muestra aumentase.
\end{enumerate}
\end{samepage}
